\clearpage\begingroup\NoSourceHead\SourceLeaf{551}
\newcommand{\BCrow}[5]{\par\noindent\hangindent=1.8em\hangafter=1 #1\enspace #2\hfill\mbox{#3}\nobreak\quad #4\enspace #5\par}
\newcommand{\BCsub}[1]{\par\smallskip\noindent #1\par}
\begin{center}\bfseries Inhaltsangabe.\end{center}
\small
\begin{center}\textls[100]{I. Teil: Die Principia generalia theoriae figurae fluidorum in statu aequilibrii.}\end{center}
\BCsub{\textit{I. Abschnitt: Überblick über die Geschichte der Extrema von Doppelintegralen vor Gauss.}}
\BCrow{1.}{Extrema von Doppelintegralen bei festem Integrationsbereich: \textsc{Lagrange} und \textsc{Euler}}{.\ .}{Seite}{4}
\BCrow{2.}{Allgemeine Charakterisierung der Methode der Variation der Grenzen und der Methode der Variation der unabhängigen Variabeln}{.\ .\ .\ .\ .\ .\ .\ .\ .\ .\ .\ .\ .\ .\ .\ .\ .\ .\ .}{—}{6}
\begin{center}A) Methode der Variation der Grenzen:\end{center}
\BCrow{3.}{Die erste \textsc{Ohm}sche Fundamentalformel}{.\ .\ .\ .\ .\ .\ .\ .\ .\ .\ .\ .\ .\ .\ .\ .\ .\ .\ .}{—}{7}
\BCrow{4.}{Die zweite \textsc{Ohm}sche Fundamentalformel}{.\ .\ .\ .\ .\ .\ .\ .\ .\ .\ .\ .\ .\ .\ .\ .\ .\ .}{—}{9}
\BCrow{5.}{Schwäche der Methode der Variation der Grenzen}{.\ .\ .\ .\ .\ .\ .\ .\ .\ .\ .\ .\ .\ .\ .}{—}{10}
\begin{center}B) Methode der Variation der unabhängigen Variabeln:\end{center}
\BCrow{6.}{Methode der Variation der unabhängigen Variabeln bei einfachen Integralen}{.\ .\ .\ .\ .}{—}{12}
\BCrow{7.}{\textsc{Euler}s erster Versuch der Ausdehnung auf Doppelintegrale}{.\ .\ .\ .\ .\ .\ .\ .\ .\ .\ .}{—}{14}
\BCrow{8.}{Die \textsc{Lagrange}sche Fundamentalformel}{.\ .\ .\ .\ .\ .\ .\ .\ .\ .\ .\ .\ .\ .\ .\ .\ .\ .\ .}{—}{16}
\BCrow{9.}{\textsc{Poisson}s Beweis derselben}{.\ .\ .\ .\ .\ .\ .\ .\ .\ .\ .\ .\ .\ .\ .\ .\ .\ .\ .\ .\ .\ .\ .}{—}{18}
\BCsub{\textit{II. Abschnitt: Die Gausssche Arbeit „Principia generalia theoriae figurae fluidorum in statu aequilibrii“.}}
\BCrow{10.}{Das \textsc{Gauss}sche Variationsproblem. Vorbemerkungen}{.\ .\ .\ .\ .\ .\ .\ .\ .\ .\ .\ .\ .}{—}{20}
\BCsub{A) Aufstellung der ersten Variation.}
\BCrow{11.}{Aufstellung der ersten Variation $\delta U$ des Flächeninhalts}{.\ .\ .\ .\ .\ .\ .\ .\ .\ .\ .\ .}{—}{23}
\BCrow{12.}{Bedeutung der \textsc{Gauss}schen Methode für die allgemeine Variationsrechnung}{.\ .\ .\ .\ .}{—}{25}
\BCrow{13.}{Berechnung von $\delta s$, $\delta\int z\,ds$, $\delta T$}{.\ .\ .\ .\ .\ .\ .\ .\ .\ .\ .\ .\ .\ .\ .\ .\ .\ .\ .\ .\ .}{—}{27}
\BCsub{B) Umformung der ersten Variation durch partielle Integration.}
\BCrow{14.}{Die \textsc{Gauss}sche Form der \textsc{Green}schen Formel}{.\ .\ .\ .\ .\ .\ .\ .\ .\ .\ .\ .\ .\ .\ .\ .}{—}{29}
\BCrow{15.}{Prioritätsfrage zwischen \textsc{Gauss} und \textsc{Green}}{.\ .\ .\ .\ .\ .\ .\ .\ .\ .\ .\ .\ .\ .\ .\ .\ .}{—}{31}
\BCrow{16.}{Die \textsc{Green}sche Formel bei \textsc{Bordoni} und \textsc{Poisson}}{.\ .\ .\ .\ .\ .\ .\ .\ .\ .\ .\ .\ .\ .}{—}{32}
\BCrow{17.}{Die \textsc{Gauss}sche Formel für $\delta W$}{.\ .\ .\ .\ .\ .\ .\ .\ .\ .\ .\ .\ .\ .\ .\ .\ .\ .\ .\ .\ .}{—}{34}
\BCsub{C) Die partielle Differentialgleichung für die Flüssigkeitsoberfläche und die Randbedingung.}
\BCrow{18.}{Ableitung der partiellen Differentialgleichung}{.\ .\ .\ .\ .\ .\ .\ .\ .\ .\ .\ .\ .\ .\ .\ .}{—}{35}
\BCrow{19.}{Ableitung der Randbedingung}{.\ .\ .\ .\ .\ .\ .\ .\ .\ .\ .\ .\ .\ .\ .\ .\ .\ .\ .\ .\ .}{—}{37}
\BCrow{20.}{Bedeutung der \textsc{Gauss}schen Resultate für die allgemeine Theorie; die Randbedingung bei \textsc{Bordoni} und \textsc{Poisson}}{.\ .\ .\ .\ .\ .\ .\ .\ .\ .\ .\ .\ .\ .\ .\ .\ .\ .\ .\ .\ .\ .\ .\ .}{—}{38}
\BCrow{21.}{Mögliche Einwände gegen die \textsc{Gauss}sche Schlussweise}{.\ .\ .\ .\ .\ .\ .\ .\ .\ .\ .\ .}{—}{40}
\BCsub{\textit{III. Abschnitt: Wirkungen der Gaussschen Arbeit:}}
\BCrow{22.}{Wirkung der \textsc{Gauss}schen Arbeit auf die allgemeine Theorie der Extrema von Doppelintegralen}{.\ .\ .\ .\ .\ .\ .\ .\ .\ .\ .\ .\ .\ .\ .\ .\ .\ .\ .\ .\ .\ .\ .\ .\ .\ .\ .\ .\ .}{—}{42}
\clearpage\SourceLeaf{552}
\ArithmeticHead{95}{R}{INHALTSANGABE.}
\BCrow{23.}{Ableitung der \textsc{Gauss}schen Resultate aus der allgemeinen Theorie: \textsc{Pagani}, \textsc{Mainardi}, \textsc{Jellett}}{.\ .\ .\ .\ .\ .\ .\ .\ .\ .\ .\ .\ .\ .\ .\ .\ .\ .\ .\ .\ .\ .\ .\ .\ .\ .\ .\ .\ .}{Seite}{44}
\BCrow{24.}{Geometrische Beweise der \textsc{Gauss}schen Resultate: \textsc{Bertrand}, \textsc{Lamarle}, \textsc{Weber}}{.\ .\ .}{—}{45}
\BCrow{25.}{Lösung des \textsc{Gauss}schen Problems in Parameterdarstellung: \textsc{Weber}, \textsc{Kneser}}{.\ .\ .\ .}{—}{46}
\BCrow{26.}{Methode der Berechnung von $\delta U$ von \textsc{Minkowski}}{.\ .\ .\ .\ .\ .\ .\ .\ .\ .\ .\ .\ .\ .}{—}{48}
\BCrow{27.}{Anwendungen der \textsc{Gauss}schen Formel für $\delta U$ auf Scharen von Minimalflächen: \textsc{Lamarle}, \textsc{Schwarz}}{.\ .\ .\ .\ .\ .\ .\ .\ .\ .\ .\ .\ .\ .\ .\ .\ .\ .\ .\ .\ .\ .\ .\ .\ .\ .\ .\ .\ .}{—}{50}
\begin{center}\textls[100]{II. Teil: Die Disquisitiones generales circa superficies curvas.}\end{center}
\BCsub{\textit{A. Die Differentialgleichung der geodätischen Linien.}}
\BCrow{a)}{Die Fläche ist durch eine Gleichung gegeben}{.\ .\ .\ .\ .\ .\ .\ .\ .\ .\ .\ .\ .\ .\ .\ .\ .}{—}{52}
\BCrow{b)}{Die Fläche ist in Parameterdarstellung gegeben}{.\ .\ .\ .\ .\ .\ .\ .\ .\ .\ .\ .\ .\ .\ .\ .}{—}{54}
\BCrow{\textit{B.}}{\textit{Die Sätze über geodätische Polar- und Parallelkoordinaten}}{.\ .\ .\ .\ .\ .\ .\ .\ .\ .\ .\ .\ .}{—}{55}
\BCsub{\hspace*{1em}1) Der Transversalensatz.}
\BCrow{\hspace*{1em}a)}{Der \textsc{Gauss}sche Satz für geodätische Linien}{.\ .\ .\ .\ .\ .\ .\ .\ .\ .\ .\ .\ .\ .\ .\ .\ .}{—}{56}
\BCrow{\hspace*{1em}b)}{Der \textsc{Kneser}sche Transversalensatz}{.\ .\ .\ .\ .\ .\ .\ .\ .\ .\ .\ .\ .\ .\ .\ .\ .\ .\ .\ .}{—}{57}
\BCrow{\hspace*{1em}c)}{Verwandte Sätze und weitere Verallgemeinerungen}{.\ .\ .\ .\ .\ .\ .\ .\ .\ .\ .\ .\ .\ .}{—}{58}
\BCrow{2)}{Beziehungen zum Unabhängigkeitssatz}{.\ .\ .\ .\ .\ .\ .\ .\ .\ .\ .\ .\ .\ .\ .\ .\ .\ .\ .}{—}{60}
\BCrow{\hspace*{1em}a)}{Die \textsc{Gauss}schen Formeln}{.\ .\ .\ .\ .\ .\ .\ .\ .\ .\ .\ .\ .\ .\ .\ .\ .\ .\ .\ .\ .\ .\ .}{—}{61}
\BCrow{\hspace*{1em}b)}{Die Untersuchungen von \textsc{Beltrami} und \textsc{Hilbert}}{.\ .\ .\ .\ .\ .\ .\ .\ .\ .\ .\ .\ .\ .}{—}{62}
\BCrow{\hspace*{1em}c)}{Verwandte Untersuchungen von \textsc{Hamilton} und weitere Verallgemeinerungen}{.\ .\ .\ .}{—}{63}
\BCrow{3)}{Der \textsc{Darboux}–\textsc{Kneser}sche Hinlänglichkeitsbeweis}{.\ .\ .\ .\ .\ .\ .\ .\ .\ .\ .\ .\ .\ .\ .}{—}{66}
\BCsub{\textit{C. Die geodätische Krümmung und der Satz von der Totalkrümmung.}}
\BCrow{1)}{Die geodätische Krümmung und ihre Verallgemeinerung}{.\ .\ .\ .\ .\ .\ .\ .\ .\ .\ .\ .\ .}{—}{70}
\BCrow{2)}{Der \textsc{Gauss}sche Satz über die Totalkrümmung und seine \textsc{Bonnet}sche Verallgemeinerung}{.}{—}{73}
\BCrow{3)}{Die \textsc{Landsberg}sche Verallgemeinerung des \textsc{Gauss}-\textsc{Bonnet}schen Satzes}{.\ .\ .\ .\ .\ .\ .}{—}{74}
\begin{center}\textls[100]{III. Teil: Die Arbeit »Allgemeine Lehrsätze in Beziehung auf die im verkehrten Verhältnisse des Quadrats der Entfernung wirkenden Anziehungs- und Abstossungskräfte.}\end{center}
\BCrow{1.}{Formulierung des \textsc{Gauss}schen Variationsproblems}{.\ .\ .\ .\ .\ .\ .\ .\ .\ .\ .\ .\ .\ .\ .\ .\ .}{—}{77}
\BCrow{2.}{A priori Beweis der Existenz einer Lösung}{.\ .\ .\ .\ .\ ,\ .\ .\ .\ .\ .\ .\ .\ .\ .\ .\ .\ .\ .}{—}{77}
\BCrow{3.}{Die Minimumsbedingung für die belegten Flächenteile}{.\ .\ .\ .\ .\ .\ ,\ .\ .\ .\ .\ .\ .\ .\ .}{—}{79}
\BCrow{4.}{Die Minimumsbedingung für die nicht belegten Flächenteile}{.\ .\ .\ .\ .\ .\ .\ .\ .\ .\ .\ .\ .}{—}{81}
\BCrow{5.}{Charakterisierung des allgemeinen Typus von Variationsproblemen, welchem das \textsc{Gauss}sche Problem angehört}{.\ .\ .\ .\ .\ .\ .\ .\ .\ .\ .\ .\ .\ .\ .\ .\ .\ .\ .\ .\ .\ .\ .\ .\ .\ .\ .\ .}{—}{82}
\BCrow{6.}{Zusammenfassung}{.\ .\ .\ .\ .\ .\ .\ .\ .\ .\ .\ .\ .\ .\ .\ .\ .\ .\ .\ .\ .\ .\ .\ .\ .\ .\ .\ .}{—}{83}
\begin{center}\textls[100]{IV. Teil: Vereinzelte kürzere Anwendungen der Variationsrechnung.}\end{center}
\BCrow{1.}{Über das Extremum des Integrals $\int n(x,y)\,ds$}{.\ .\ .\ .\ .\ .\ .\ .\ .\ .\ .\ .\ .\ .\ .\quad .\ .}{—}{84}
\BCrow{2.}{Über die elastische Kurve}{.\ .\ .\ .\ .\ .\ .\ .\ .\ .\ .\ .\ .\ .\ .\ .\ .\ .\ .\ .\ .\ .\ .\ .\ .}{—}{85}
\BCrow{3.}{Über den homogenen Körper von gegebener Masse und grösster Anziehung}{.\ .\ .\ .\ .\ .\ .}{—}{88}
\begin{center}\rule{3.5cm}{.25pt}\end{center}
\clearpage\endgroup
